% Created 2026-05-11 Mon 02:20
% Intended LaTeX compiler: lualatex
\documentclass[11pt]{article}

\usepackage{amsmath}
\usepackage{fontspec}
\usepackage{graphicx}
\usepackage{longtable}
\usepackage{wrapfig}
\usepackage{rotating}
\usepackage[normalem]{ulem}
\usepackage{capt-of}
\usepackage{hyperref}
\usepackage{minted}
\usepackage{fontspec}
\setmainfont{UT Sans}[
BoldFont={UT Sans Bold},
UprightFont={UT Sans Regular},
FontFace={m}{n}{UT Sans Medium}
]
% Fallback fake italic using a slant
\usepackage{etoolbox}
\newcommand{\fakeitalic}[1]{{\addfontfeatures{FakeSlant=0.18}#1}}
\AtBeginDocument{\renewcommand{\emph}[1]{\fakeitalic{#1}}}
\usepackage{minted}
\usepackage[a4paper,margin=3cm]{geometry}
\usepackage{lastpage}
\usepackage{fancyhdr}
\usepackage{lastpage}
\pagestyle{fancy}                % set fancy style
\fancyhf{}                        % clear all headers and footers
\cfoot{\thepage/\pageref*{LastPage}}  % center footer
\renewcommand{\headrulewidth}{0pt}    % remove header line
\renewcommand{\footrulewidth}{0pt}    % remove footer line
% also override plain page style for first page
\fancypagestyle{plain}{%
\fancyhf{}%
\cfoot{\thepage/\pageref*{LastPage}}%
\renewcommand{\headrulewidth}{0pt}%
\renewcommand{\footrulewidth}{0pt}%
}
\usepackage{url}
\author{Elias-Robert BAJCSI}
\date{}
\title{Laborator 4\\\medskip
\large PLFP}
\hypersetup{
 pdfauthor={Elias-Robert BAJCSI},
 pdftitle={Laborator 4},
 pdfkeywords={},
 pdfsubject={},
 pdfcreator={},
 pdflang={English}}
\begin{document}

\maketitle
\section{\texttt{./Cargo.toml}}
\label{sec:org6fb5f2d}
\begin{minted}[breaklines=true,breakanywhere=true,breaksymbol=,fontsize=\small,linenos=false]{toml}
[package]
name = "lab4"
version = "0.1.0"
edition = "2021"

[[bin]]
name = "ex1"
path = "src/bin/ex1.rs"
[[bin]]
name = "ex2"
path = "src/bin/ex2.rs"
[[bin]]
name = "ex3"
path = "src/bin/ex3.rs"
[[bin]]
name = "ex4"
path = "src/bin/ex4.rs"
\end{minted}
\section{\texttt{./src/bin/ex1.rs}}
\label{sec:org20c6d9f}
\begin{minted}[breaklines=true,breakanywhere=true,breaksymbol=,fontsize=\small,linenos=false]{rs}
use std::cmp::Ordering;

#[derive(Debug, Clone)]
struct Album {
    titlu: String,
    gen: String,
}

fn cel_mai_bun_dupa<'a>(albume: &'a [Album], cheie: fn(&Album) -> &str) -> Option<&'a Album> {
    albume.iter().min_by(|a, b| cheie(a).cmp(cheie(b)))
}

fn main() {
    let mut albume = vec![
        Album { titlu: "Black Monday".into(), gen: "Metal".into() },
        Album { titlu: "Scientist".into(), gen: "Dub".into() },
        Album { titlu: "Fireworks".into(), gen: "Pop".into() },
        Album { titlu: "Fade".into(), gen: "Pop".into() },
        Album { titlu: "Undertow".into(), gen: "Metal".into() },
    ];

    albume.sort_by(|a, b| {
        let g = a.gen.cmp(&b.gen);
        if g == Ordering::Equal { a.titlu.cmp(&b.titlu) } else { g }
    });

    for a in &albume {
        println!("{} - {}", a.gen, a.titlu);
    }

    let minim_dupa_titlu = cel_mai_bun_dupa(&albume, |a| &a.titlu).unwrap();
    println!("min titlu: {:?}", minim_dupa_titlu);

    let div7: Vec<i32> = (1..=1000).filter(|n| n % 7 == 0).collect();
    let suma: i32 = div7.iter().sum();
    println!("count: {}, suma: {}", div7.len(), suma);
}
\end{minted}
\section{\texttt{./src/bin/ex2.rs}}
\label{sec:org80a1376}
\begin{minted}[breaklines=true,breakanywhere=true,breaksymbol=,fontsize=\small,linenos=false]{rs}
fn make_squares_up_to(x: u64) -> Vec<u64> {
    (0..x).map(|n| n * n).collect()
}

fn factorial(n: u64) -> u64 {
    if n <= 1 { 1 } else { n * factorial(n - 1) }
}

fn make_factorial() -> impl FnMut(u64) -> u64 {
    move |n| factorial(n)
}

fn make_running_sum() -> impl FnMut(i32) -> i32 {
    let mut suma = 0;
    move |x| {
        suma += x;
        suma
    }
}

fn main() {
    println!("{:?}", make_squares_up_to(10));

    let mut fact = make_factorial();
    println!("{} {} {}", fact(3), fact(5), fact(7));

    let mut run = make_running_sum();
    println!("{} {} {}", run(3), run(5), run(2));
}
\end{minted}
\section{\texttt{./src/bin/ex3.rs}}
\label{sec:orgb30c6de}
\begin{minted}[breaklines=true,breakanywhere=true,breaksymbol=,fontsize=\small,linenos=false]{rs}
fn agregeaza<T, R>(colectie: &[T], f: impl Fn(&[T]) -> R) -> R {
    f(colectie)
}

fn mediana(nums: &[i32]) -> Option<i32> {
    if nums.is_empty() {
        return None;
    }
    let mut v = nums.to_vec();
    v.sort();
    Some(v[v.len() / 2])
}

fn main() {
    let nums = [9, 3, 1, 8, 7, 2, 6, 5, 4];

    let suma = agregeaza(&nums, |v| v.iter().sum::<i32>());
    let maxim = agregeaza(&nums, |v| *v.iter().max().unwrap());
    let med = agregeaza(&nums, mediana);

    println!("suma: {}", suma);
    println!("maxim: {}", maxim);
    println!("mediana: {:?}", med);
}
\end{minted}
\section{\texttt{./src/bin/ex4.rs}}
\label{sec:org35cdfd1}
\begin{minted}[breaklines=true,breakanywhere=true,breaksymbol=,fontsize=\small,linenos=false]{rs}
fn compune_n(fns: Vec<Box<dyn Fn(i32) -> i32>>) -> impl Fn(i32) -> i32 {
    move |x| fns.iter().fold(x, |acc, f| f(acc))
}

fn main() {
    let fns: Vec<Box<dyn Fn(i32) -> i32>> = vec![
        Box::new(|x| x + 1),
        Box::new(|x| x * 2),
        Box::new(|x| x - 3),
        Box::new(|x| x * x),
    ];

    let compusa = compune_n(fns);
    println!("{}", compusa(5));

    println!(
        "La [Album], agregarea trebuie sa foloseasca un camp numeric (ex: durata, rating), nu varste inexistente."
    );
}
\end{minted}
\end{document}
